\SourcePage{540}{543}
\section{Matrix elements in the addition of angular momenta}

Consider once more a system made up of two parts (to be called subsystems 1 and 2), and let $f_{kq}^{(1)}$ be a spherical tensor characterizing the first subsystem. According to (107.6), its matrix elements with respect to the wave functions of that same subsystem are defined by
\begin{equation}
 \begin{aligned}
 \langle n'_1j'_1m'_1|f_{kq}^{(1)}|n_1j_1m_1\rangle
 ={}&i^k(-1)^{j_{1\max}-m'_1}
 \begin{pmatrix}j'_1&k&j_1\\-m'_1&q&m_1\end{pmatrix}\
 &\times\langle n'_1j'_1\|f_k^{(1)}\|n_1j_1\rangle.
 \end{aligned}
 \tag{109.1}
\end{equation}

\SourcePage{541}{544}
We now ask how to calculate matrix elements of these same quantities with respect to wave functions of the complete system. We shall show how they can be expressed in terms of the same reduced matrix elements that occur in (109.1).

The states of the complete system are specified by the quantum numbers $j_1j_2JMn_1n_2$ ($J$ and $M$ are the magnitude and projection of the angular momentum of the complete system). Since $f_{kq}^{(1)}$ belongs to subsystem 1, its operator commutes with the angular-momentum operator of subsystem 2. Its matrix is consequently diagonal in $j_2$, and also in the other quantum numbers $n_2$ of that subsystem. For brevity we shall omit these indices ($j_2$, $n_2$) and write the matrix elements sought as
\[
 \langle n'_1j'_1J'M'|f_{kq}^{(1)}|n_1j_1JM\rangle.
\]
By (107.6), their dependence on $M$ is given by
\begin{equation}
 \begin{aligned}
 \langle n'_1j'_1J'M'|f_{kq}^{(1)}|n_1j_1JM\rangle
 ={}&i^k(-1)^{J_{\max}-M'}
 \begin{pmatrix}J'&k&J\\-M'&q&M\end{pmatrix}\
 &\times\langle n'_1j'_1J'\|f_k^{(1)}\|n_1j_1J\rangle.
 \end{aligned}
 \tag{109.2}
\end{equation}

To establish the relation between the reduced matrix elements on the right-hand sides of (109.1) and (109.2), we write, by the definition of matrix elements,
\[
 \begin{aligned}
 \langle n'_1j'_1J'M'|f_{kq}^{(1)}|n_1j_1JM\rangle
 & =\int\psi_{J'M'}^{*}\widehat f_{kq}^{(1)}\psi_{JM}\,dq={}\\
 & =\sum_{m_1m'_1}(-1)^{j'_1-j_2+M'-M}
 \sqrt{(2J'+1)(2J+1)}
 \begin{pmatrix}j'_1&j_2&J'\\m'_1&m_2&-M'\end{pmatrix}\\
 &\qquad\times
 \begin{pmatrix}j_1&j_2&J\\m_1&m_2&-M\end{pmatrix}
 \langle n'_1j'_1m'_1|f_{kq}^{(1)}|n_1j_1m_1\rangle.
 \end{aligned}
\]
Substitution of (109.1) and (109.2) into this expression, followed by comparison of the resulting relation with (108.4), shows that the ratio of the reduced matrix elements in (109.1) and (109.2) must be proportional to a particular $6j$ symbol. Careful comparison of the two relations gives the follo\SourcePageJoin{542}{545}wing final formula:
\begin{equation}
 \begin{aligned}
 \langle n'_1j'_1J'\|f_k^{(1)}\|n_1j_1J\rangle
 ={}&(-1)^{j_{1\max}+j_2+J_{\min}+k}
 \sqrt{(2J+1)(2J'+1)}\\
 &\times
 \begin{Bmatrix}j'_1&J'&j_2\\J&j_1&k\end{Bmatrix}
 \langle n'_1j'_1\|f_k^{(1)}\|n_1j_1\rangle
 \end{aligned}
 \tag{109.3}
\end{equation}
(here $j_{1\max}$ is the larger of $j_1$, $j'_1$, while $J_{\min}$ is the smaller of $J$, $J'$). The corresponding formula for reduced matrix elements of a spherical tensor belonging to the second subsystem is
\begin{equation}
 \begin{aligned}
 \langle n'_2j'_2J'\|f_k^{(2)}\|n_2j_2J\rangle
 ={}&(-1)^{j_1+j_{2\min}+J_{\max}+k}
 \sqrt{(2J+1)(2J'+1)}\\
 &\times
 \begin{Bmatrix}j'_2&J'&j_1\\J&j_2&k\end{Bmatrix}
 \langle n'_2j'_2\|f_k^{(2)}\|n_2j_2\rangle.
 \end{aligned}
 \tag{109.4}
\end{equation}

The lack of complete symmetry between (109.3) and (109.4), in the exponent of $-1$, is connected with the dependence of the wave-function phase on the order in which the angular momenta are added. This difference must be kept in mind when matrix elements have to be evaluated for both subsystems at once.

We next derive a useful formula for matrix elements, with respect to wave functions of the complete system, of the scalar product (see definition (107.4)) of two spherical tensors of the same rank $k$ which belong to different subsystems (and therefore commute with one another). According to (107.10), these matrix elements are expressed through the reduced matrix elements of each tensor (with respect to wave functions of the complete system) as follows:
\[
 \begin{aligned}
 &\langle n'_1n'_2j'_1j'_2JM|(f_k^{(1)}f_k^{(2)})_{00}|n_1n_2j_1j_2JM\rangle={}\\
 &\quad=\frac1{2J+1}\sum_{J''}
 \langle n'_1j'_1J\|f_k^{(1)}\|n_1j_1J''\rangle
 \langle n'_2j'_2J''\|f_k^{(2)}\|n_2j_2J\rangle
 \end{aligned}
\]
(we have used here the fact that the matrix of a quantity belonging to either subsystem is diagonal in the quantum numbers of the other subsystem). Substituting (109.3) and (109.4) here and using the summation formula (108.8), we obtain the desired formula, which expresses the matrix elements of the scalar product through the reduced matrix elements of each tensor with respect to the wave functions of the correspo\SourcePageJoin{543}{546}nding subsystems:
\begin{equation}
 \begin{aligned}
 &\langle n'_1n'_2j'_1j'_2JM|(f_k^{(1)}f_k^{(2)})_{00}|n_1n_2j_1j_2JM\rangle={}\\
 &\quad=(-1)^{j_{1\min}+j_{2\max}+J}
 \begin{Bmatrix}J&j'_2&j'_1\\k&j_1&j_2\end{Bmatrix}\\
 &\qquad\times
 \langle n'_1j'_1\|f_k^{(1)}\|n_1j_1\rangle
 \langle n'_2j'_2\|f_k^{(2)}\|n_2j_2\rangle.
 \end{aligned}
 \tag{109.5}
\end{equation}
